\scanpage{38}
\uline{Proposition} \uline{61} \inserted{$\forall n\in\mathbb N$} L'hom $j_n$ est injectif, et on
a commutativité dans le diagramme
\diagram{\begin{tikzcd}[ampersand replacement=\&,column sep=small,row sep=large]
L^{\otimes -n}\otimes\underline{\mathrm{Sym}}^{\unsure{n|*}}(V)\arrow[r,"\alpha_n","\sim"']\arrow[d,"i_n","{\mathrm{id}_{L^{\otimes(-n-1)}}\otimes\deleted{\unsure{i_n(L)}}}"'] \& \underline{\mathrm{Sym}}^n(\check V)\arrow[r,hook,"j_n"]\arrow[d,hook,"{\substack{\text{multiplication par}\\e_{\check V}=i'_1(1)}}"] \& \underline{\mathrm{Alg\,aff}}(F)\arrow[d,equal]\\
L^{\otimes(-n-1)}\otimes\mathrm{Sym}^{n+1}(V)\arrow[r,"\sim","\alpha_{n+1}"'] \& \underline{\mathrm{Sym}}^{n+1}(\check V)\arrow[r,hook,no head,"j_{n+1}"] \& \mathrm{Alg\,aff}(F).
\end{tikzcd}}
\margin{\deleted{\unsure{$i$}}}
Enfin, $\mathrm{Alg\,aff}(F)$ est isomorphe à la limite
inductive des $\mathrm{Sym}^n(\check V)\simeq L^{\otimes(-n)}\otimes\mathrm{Sym}^n(V)$.

\uline{NB} Les sections de $\mathrm{Sym}^n(\check V)$, identifiées à des
fonctions $\unsure{f|(-}$ sur $F$, s'appellent les fonctions \inserted{polynômiales} de
degré $\leq n$. Correspondance entre celles-ci et
formes sur \inserted{$V($\unsure{$=$}}$\check{\mathbb V}(V))$ de degré $n$. Pour \deleted{\unsure{tt|une}} telle
$f$, sa classe modulo $\mathrm{Sym}^{n-1}(\check V)$ (fonctions
polynômiales de degré $\leq n-1$) s'appelle
\uline{terme dominant} de $f$, et s'identifie
à une section de $L^{\otimes(-n)}$; \deleted{\unsure{c'est|ainsi}} \inserted{elle s'identifie} aussi à
la restriction à $L$ de la $n$-forme sur $V$ associée
à $f$. --- On notera que dans cette proposition
\margin{$!\,\left\{\vphantom{\begin{matrix}1\\1\\1\end{matrix}}\right.$}
se trouvent mélangés des \unsure{résultats} \inserted{valables} sur les
fibrés affines de rang quelconque, et ceux
spéciaux aux fibrés de rang 1, liés à
la dualité (52).
% archive sheet number: 37
